Tables are among the most pervasive means of organizing information. From relational databases that power enterprise systems, to spreadsheets used in finance and science, to the info-boxes and result tables embedded in web pages and research papers, a large fraction of the world's structured knowledge is stored in tabular form. Enabling machines to \emph{understand} this data (to answer questions over it, verify claims against it, summarize it, and translate natural-language intent into precise queries against it) is therefore a foundational problem for natural language processing and a prerequisite for democratizing access to data. The stakes are highest in the enterprise: organizations sit on vast operational databases and warehouses, yet the ability to interrogate them is gated behind specialists who can write SQL and who know what each cryptic column actually means. Letting any analyst or decision-maker ask a question in natural language and get a trustworthy answer over data at this scale is the practical prize that motivates the field, and the gap between that promise and what current systems deliver is the problem this report sets out to frame.

The recent progress of Large Language Models (LLMs) and Multimodal Large Language Models (MLLMs) has renewed interest in tabular data, but tables expose a fundamental mismatch with how these models process input. LLMs are trained on \emph{linear}, one-dimensional token sequences, whereas a table is inherently \emph{two-dimensional}: its meaning is carried jointly by rows, columns, headers, and the relationships between cells~\cite{wu2025tabular}. Compounding this, tables are structurally \emph{flexible}. They range from well-normalized relational database tables, to hierarchical and multi-layered spreadsheets, to semi-structured info-boxes linked to free text. This diversity of format and purpose has driven the field toward a fragmented landscape of specialized representations, tasks, and datasets rather than universal solutions, making the area difficult to navigate and models difficult to generalize~\cite{wu2025tabular}.

Before any reasoning can occur, a two-dimensional table must be converted into a form a model can consume. The choice of this input representation is not neutral: recent studies show that model performance can swing substantially depending on how closely the chosen format aligns with the data seen during pre-training~\cite{sui2024table}. The literature broadly distinguishes four families of representation: \emph{serialization} (e.g., HTML, JSON, Markdown, or LaTeX renderings of the full table), \emph{data schema} (the structural blueprint of the data, such as SQL \texttt{CREATE} statements or pandas dataframes), \emph{image}-based formats, and specialized \emph{table encoders}~\cite{wu2025tabular}. These representations trade off completeness against compactness: serialization and images preserve content but are bounded by context length and resolution, whereas a schema-based view bypasses length limits by exposing only structure, at the cost of the concrete values. A limitation that is easy to overlook is one of scale: these representations, and the methods built on them, are overwhelmingly validated on \emph{small} tables, often no larger than four or five columns and a few dozen rows. Real-world tables are nothing like this. Enterprise databases hold hundreds of columns and thousands to millions of rows, at which point serialization overflows the context window, image renderings blur past legibility, and even schema-only views strain as the number of tables and relationships explodes. A further, more structural option casts the table as a \emph{graph}, making its layout explicit and addressable rather than implicit in a flattened token stream; this is a promising foundation, but current graph methods are demonstrated only on small or document-centric data, and extending them to large tables is part of the motivation this report develops. The chasm between the toy tables on which solutions are demonstrated and the industry-scale tables on which they must actually work is the central motivation of this report.

This report studies tabular data understanding through two complementary lenses that together span this space:

\begin{itemize}
    \item \textbf{Schema understanding} is concerned with the \emph{structural blueprint} of tabular data (table and column names, data types, primary and foreign keys, and inter-table relationships) and with mapping a user's natural-language intent onto that structure. The difficulty is defined by industry reality: enterprise schemas are large (thousands of tables and tens of thousands of columns), their naming is cryptic and undocumented (e.g., columns named \texttt{A2} or \texttt{dname}), and they are \emph{obsolescent}, drifting as columns are added, repurposed, or abandoned while documentation lags behind. The central challenges are therefore disambiguating obscure and stale schemas, retrieving the relevant schema elements for a query from among tens of thousands of candidates, and generating structurally and semantically valid operations such as SQL over them. Text-to-SQL is the canonical task in this pillar.

    \item \textbf{Data understanding} is concerned with reasoning over the \emph{content} of tables (the values, their aggregation, and the trends, outliers, and relationships they encode) often in combination with surrounding text. Here the challenges are higher-order reasoning that goes beyond retrieval and arithmetic, and robustness to complex inputs such as hierarchical, long, and multi-table settings. The same scale problem recurs: methods excel on the small, clean tables of academic benchmarks but degrade sharply once a table is large enough that the handful of answer-bearing cells is buried among thousands of irrelevant ones. Table Question Answering is the canonical task in this pillar, with table fact verification and table-to-text generation as adjacent tasks.
\end{itemize}

These two pillars are complementary rather than disjoint. A schema is a lossy abstraction of the underlying data, while the data is only interpretable in light of its schema; robust tabular understanding requires both. The current state of the field, however, is uneven, and in a way that sharpens the motivation above. On the schema-understanding side, Text-to-SQL has advanced rapidly from early encoder--decoder architectures~\cite{zhong2017seq2sql} to LLM-based in-context and multi-agent pipelines~\cite{pourreza2023din,wang2024mac,pourreza2024chase}, with strong results on benchmarks such as BIRD~\cite{li2023bird}, yet these gains are reported on curated databases and do not transfer cleanly to the obscure, obsolescent, industry-scale schemas described above. On the data-understanding side, several benchmarks are approaching saturation on retrieval-oriented questions over small tables, yet models remain markedly weaker than humans once inputs become structurally complex: on hierarchical and multi-table benchmarks such as HiTab~\cite{cheng2022hitab} and MULTIHIERTT~\cite{zhao2022multihiertt}, reported accuracies fall well below human performance. The unifying motivation of this report is therefore to close the gap between demonstration and deployment: to make \emph{both} schema understanding and data understanding work at industry scale, where schemas are large and stale and tables are far too big to fit in a model's context. Schema understanding is the natural first front, and data understanding the next frontier once it is well founded.

%----------------------------------------------
\section{Objective and scope}
The objective of this report is to survey and consolidate current research on tabular data understanding with large language models, organized around the two pillars of schema understanding and data understanding, and to situate our own research within this landscape. The report is scoped as follows:

\begin{itemize}
    \item We treat the taxonomy of tabular input representations (serialization, data schema, image, table encoders, and graph-based representations) within the data-understanding pillar, since this choice bears most directly on reasoning over content, and examine the sensitivity of model behaviour to it.

    \item Under \textbf{schema understanding}, we review knowledge-base and metadata construction for obscure schemas, schema linking and retrieval, and the evolution of Text-to-SQL methods, culminating in retrieval-augmented and self-reflective LLM pipelines.

    \item Under \textbf{data understanding}, we center on Table Question Answering, noting table fact verification and table-to-text as adjacent tasks, and review the open challenges of reasoning beyond retrieval, robustness to complex inputs, and generalization across tabular representations.

    \item As current progress, we present \emph{Reflect-SQL}, a self-reflection-based framework for Text-to-SQL that operationalizes the schema-understanding pillar, and we use it to motivate data understanding as the forward research direction.

    \item The report focuses on methods, tasks, and benchmarks at the level needed to frame a research programme; it is not intended as an exhaustive technical account of every approach.
\end{itemize}

%----------------------------------------------
\section{Contributions}
The following are the contributions of this report:
\begin{itemize}
    \item A structured survey of tabular data understanding that unifies the fragmented literature under two complementary pillars, schema understanding and data understanding, treating the question of how tables are represented to models (including graph-based representations) within the data-understanding pillar.

    \item A synthesis of the Text-to-SQL and schema-understanding literature, tracing the progression from encoder--decoder models to knowledge-base-driven, retrieval-augmented, and self-reflective LLM frameworks, and identifying where reliability gaps remain.

    \item A synthesis of the data-understanding literature that consolidates the principal reasoning tasks and articulates three open challenges: limited reasoning beyond retrieval, lack of robustness on complex inputs, and limited cross-representation generalization.

    \item \emph{Reflect-SQL}~\cite{jain2027reflectsql}, our framework for enterprise Text-to-SQL, which introduces a dynamically enriched knowledge base, a feedback-driven retrieval loop, an iterative SQL synthesis loop, and a reference-free entailment check, and which achieves an execution accuracy of $72.03\%$ on the BIRD benchmark.

    \item An extensive research agenda spanning both pillars: for schema understanding, continually evolving knowledge bases, retrieval that scales by addressing pre-computed structure rather than searching at query time, cross-database join reasoning, and reliable yet inexpensive reflection; and for data understanding, representing tabular data at scale as a graph as the lead direction, with evidence retrieval recast as addressing over that structure and a large-scale Table Question Answering benchmark to measure progress.
\end{itemize}

%---------------------------------------------------------
\section{Structure of the report}
The remainder of this report is organized into four chapters. Chapter~\ref{ch:ch1} presents the literature survey: organized around the two pillars, it reviews the schema-understanding pillar (knowledge-base construction, schema linking and retrieval, and the evolution of Text-to-SQL from encoder--decoder models to self-reflective LLM pipelines) and the data-understanding pillar (how table content is represented, including graph-based representations, Table Question Answering, and its open challenges), establishing the research context for both pillars in a balanced manner. Chapter~\ref{ch:ch3} reports current progress: it details \emph{Reflect-SQL}, our self-reflection framework for Text-to-SQL, including its knowledge base, retrieval and generation loops, reference-free entailment check, and experimental results on the BIRD benchmark. Chapter~\ref{ch:future} sets out an extensive research agenda: it consolidates what schema understanding has achieved and the problems that remain, then charts our next direction in data understanding, centered on representing tabular data at scale (concretely as a graph) and the large-scale benchmark needed to measure it. Chapter~\ref{ch:conc} concludes by summarizing the findings and restating the report's central thesis: that progress in tabular data understanding depends less on ever-larger models than on the architecture that supplies them with the right structure and the right evidence.
